%============================================================
%  Chapter 5 – Database Design
%============================================================
\chapter{Database Design}

\section{Overview}

\projectname{} uses \textbf{MongoDB} as the underlying database, accessed through the \textbf{Prisma ORM} (version 6.x). MongoDB's document model was chosen for its schema flexibility, horizontal scalability, and native JSON compatibility with Node.js. Prisma provides a type-safe, auto-generated client and a declarative schema language that makes the data model self-documenting.

The complete schema is defined in \texttt{prisma/schema.prisma} and comprises approximately \textbf{25 models} and \textbf{12 enumerations}. All primary keys are MongoDB ObjectIds mapped to Prisma's \texttt{@db.ObjectId} type.

\section{Enumeration Types}

The schema defines the following enumerations to enforce valid state values across models:

\begin{table}[H]
\centering
\caption{Database Enumeration Types}
\begin{tabularx}{\textwidth}{lX}
\toprule
\textbf{Enum} & \textbf{Values} \\
\midrule
\texttt{Role}               & \texttt{CLIENT}, \texttt{FREELANCER}, \texttt{ADMIN} \\
\texttt{ProjectStatus}      & \texttt{DRAFT}, \texttt{OPEN}, \texttt{IN\_PROGRESS}, \texttt{UNDER\_REVIEW}, \texttt{REVISION\_REQUESTED}, \texttt{COMPLETED}, \texttt{CANCELLED}, \texttt{DISPUTED} \\
\texttt{ProjectSize}        & \texttt{SMALL}, \texttt{MEDIUM}, \texttt{LARGE} \\
\texttt{ProposalStatus}     & \texttt{PENDING}, \texttt{SHORTLISTED}, \texttt{ACCEPTED}, \texttt{REJECTED}, \texttt{WITHDRAWN} \\
\texttt{ContractStatus}     & \texttt{ACTIVE}, \texttt{COMPLETED}, \texttt{CANCELLED}, \texttt{DISPUTED} \\
\texttt{MilestoneStatus}    & \texttt{PENDING}, \texttt{IN\_PROGRESS}, \texttt{SUBMITTED}, \texttt{APPROVED}, \texttt{REJECTED} \\
\texttt{PaymentStatus}      & \texttt{PENDING}, \texttt{HELD\_IN\_ESCROW}, \texttt{RELEASED}, \texttt{REFUNDED}, \texttt{FAILED} \\
\texttt{DisputeStatus}      & \texttt{OPEN}, \texttt{UNDER\_REVIEW}, \texttt{RESOLVED}, \texttt{CLOSED} \\
\texttt{DisputeResolution}  & \texttt{FAVOR\_CLIENT}, \texttt{FAVOR\_FREELANCER}, \texttt{SPLIT}, \texttt{DISMISSED} \\
\texttt{TokenWalletTxType}  & \texttt{CREDIT}, \texttt{DEBIT} \\
\texttt{TokenTxReason}      & \texttt{REGISTRATION\_BONUS}, \texttt{MONTHLY\_AWARD}, \texttt{PROPOSAL\_SUBMITTED}, \texttt{PROPOSAL\_WITHDRAWN}, \texttt{ADMIN\_GRANT}, \texttt{ADMIN\_DEDUCT} \\
\texttt{ExperienceLevel}    & \texttt{ENTRY}, \texttt{MID}, \texttt{SENIOR}, \texttt{EXPERT} \\
\texttt{AvailabilityStatus} & \texttt{AVAILABLE}, \texttt{BUSY}, \texttt{UNAVAILABLE} \\
\texttt{InvitationStatus}   & \texttt{PENDING}, \texttt{ACCEPTED}, \texttt{REJECTED}, \texttt{EXPIRED} \\
\texttt{MessageType}        & \texttt{TEXT}, \texttt{FILE}, \texttt{SYSTEM} \\
\texttt{NotificationType}   & \texttt{PROPOSAL\_RECEIVED}, \texttt{PROPOSAL\_ACCEPTED}, \texttt{PAYMENT\_RELEASED}, \texttt{TOKEN\_CREDITED}, and others (16 total) \\
\bottomrule
\end{tabularx}
\end{table}

\section{Core Data Models}

\subsection{User}

The \texttt{User} model is the unified authentication entity. A single \texttt{users} collection stores all account types; the \texttt{role} field determines whether the user is a Client, Freelancer, or Admin, and controls which sub-profile document is created.

\begin{table}[H]
\centering
\caption{User Model Fields}
\begin{tabularx}{\textwidth}{llX}
\toprule
\textbf{Field} & \textbf{Type} & \textbf{Description} \\
\midrule
\texttt{id}               & ObjectId & Primary key. \\
\texttt{name}             & String   & Display name. \\
\texttt{email}            & String   & Unique email address. \\
\texttt{passwordHash}     & String?  & Bcrypt hash; null for OAuth-only accounts. \\
\texttt{googleId}         & String?  & Google OAuth sub (subject) identifier. \\
\texttt{role}             & Role?    & CLIENT, FREELANCER, or ADMIN. \\
\texttt{isEmailVerified}  & Boolean  & True after OTP verification. \\
\texttt{isIdVerified}     & Boolean  & True after admin verifies the uploaded ID. \\
\texttt{idDocumentUrl}    & String?  & Cloudinary URL of the uploaded identity document. \\
\texttt{isBanned}         & Boolean  & Set by admin; blocks login. \\
\texttt{profileImage}     & String?  & Cloudinary URL of avatar. \\
\texttt{createdAt}        & DateTime & Account creation timestamp. \\
\texttt{lastActiveAt}     & DateTime? & Used by the browse scoring algorithm. \\
\bottomrule
\end{tabularx}
\end{table}

\subsection{Session}

The \texttt{Session} model stores issued refresh tokens, enabling server-side revocation and preventing token reuse attacks.

\subsection{ClientProfile and FreelancerProfile}

Each role has a dedicated profile document (stored in \texttt{client\_profiles} and \texttt{freelancer\_profiles} collections) linked via a \texttt{userId} foreign reference. This one-to-one relationship keeps the core \texttt{User} document lean while allowing role-specific fields.

Key \texttt{FreelancerProfile} fields include:

\begin{itemize}
  \item \texttt{skillTokenBalance} — current SkillToken balance.
  \item \texttt{profileCompletion} — integer 0–100 representing onboarding completeness.
  \item \texttt{experienceLevel} — \texttt{ExperienceLevel} enum.
  \item \texttt{availability} — \texttt{AvailabilityStatus} enum.
  \item \texttt{preferredCategories} — array of category slugs for browse personalisation.
  \item \texttt{preferredBudgetMin / Max} — browse filter defaults.
\end{itemize}

\subsection{Project}

The \texttt{Project} model represents a client's work posting. Key design decisions include:

\begin{itemize}
  \item \texttt{status} follows the \texttt{ProjectStatus} state machine from \texttt{DRAFT} through \texttt{COMPLETED} or \texttt{CANCELLED}.
  \item \texttt{budgetType} can be \texttt{"fixed"} or \texttt{"hourly"}.
  \item \texttt{isAiScoped} and \texttt{aiScopeSummary} record whether the OpenAI assistant was used to generate the project spec.
  \item \texttt{attachments} is a \texttt{String[]} array of Cloudinary URLs.
\end{itemize}

\subsection{Proposal}

The \texttt{Proposal} model captures a freelancer's bid on a project. The \texttt{tokenCost} field records the SkillTokens deducted at submission. A unique constraint on \texttt{[projectId, freelancerProfileId]} enforces one proposal per freelancer per project.

\subsection{Contract, Milestone, and Payment}

Once a proposal is accepted, a \texttt{Contract} document is created. Each contract has an array of \texttt{Milestone} sub-documents representing deliverable checkpoints. Each milestone has an associated \texttt{Payment} document tracking the escrow and release status. The relationship chain is:

\begin{center}
\texttt{Project} $\rightarrow$ \texttt{Contract} $\rightarrow$ \texttt{Milestone[]} $\rightarrow$ \texttt{Payment}
\end{center}

\subsection{Dispute}

The \texttt{Dispute} model links a project to the client, freelancer, and optionally an admin. The \texttt{resolution} field uses the \texttt{DisputeResolution} enum to record the outcome (\texttt{FAVOR\_CLIENT}, \texttt{FAVOR\_FREELANCER}, \texttt{SPLIT}, or \texttt{DISMISSED}).

\subsection{Review}

The \texttt{Review} model implements the blind review system. The \texttt{isRevealed} boolean (default \texttt{false}) controls visibility. Reviews are revealed when both parties have submitted or a timeout elapses. A unique constraint on \texttt{[contractId, giverId]} prevents duplicate reviews.

\subsection{TokenTransaction}

An immutable ledger of all SkillToken movements. Every row stores the \texttt{amount}, \texttt{type} (CREDIT/DEBIT), \texttt{reason} (from \texttt{TokenTxReason} enum), and a \texttt{balanceAfter} snapshot. The snapshot enables full balance reconstruction and audit trail.

\subsection{Additional Models}

\begin{table}[H]
\centering
\caption{Additional Database Models}
\begin{tabularx}{\textwidth}{lX}
\toprule
\textbf{Model} & \textbf{Purpose} \\
\midrule
\texttt{Skill}            & Global skills catalogue (name, category). \\
\texttt{FreelancerSkill}  & Many-to-many join: freelancer $\leftrightarrow$ skill (with proficiency 1–5). \\
\texttt{ProjectSkill}     & Many-to-many join: project $\leftrightarrow$ required skill. \\
\texttt{PortfolioItem}    & Freelancer's past work showcase items. \\
\texttt{Certificate}      & Freelancer's professional certificates. \\
\texttt{Education}        & Freelancer's education history. \\
\texttt{Gig}              & Freelancer's pre-packaged service offering. \\
\texttt{Category}         & Project category (e.g., Web Development). \\
\texttt{SubCategory}      & Project sub-category (e.g., Frontend). \\
\texttt{Invitation}       & Client-to-freelancer project invitation. \\
\texttt{ChatRoom}         & One chat room per contract or pre-contract. \\
\texttt{Message}          & Individual message within a chat room. \\
\texttt{Notification}     & In-app notification for any event. \\
\texttt{AdminLog}         & Immutable audit trail of admin actions. \\
\texttt{PlatformSetting}  & Key-value store for global platform config. \\
\texttt{PasswordResetToken} & Secure time-limited token for password resets. \\
\texttt{AiScopingSession} & Stored AI assistant conversation for each scoping interaction. \\
\texttt{SavedProject}     & Freelancer bookmarks for projects. \\
\texttt{BrowseInteraction} & Tracks freelancer project views for personalisation. \\
\bottomrule
\end{tabularx}
\end{table}

\section{Database Design Principles}

\begin{enumerate}
  \item \textbf{Cascade deletes}: All child documents use \texttt{onDelete: Cascade} so that deleting a user removes all associated profiles, projects, proposals, and messages automatically.
  \item \textbf{Unique constraints}: Applied on email (User), refresh tokens (Session), and join-table composite keys to prevent data duplication.
  \item \textbf{Immutability for ledgers}: \texttt{TokenTransaction} and \texttt{AdminLog} records are never updated or deleted, ensuring a full audit trail.
  \item \textbf{Snapshot fields}: \texttt{TokenTransaction.balanceAfter} stores the running balance at transaction time, enabling O(1) current-balance lookup.
  \item \textbf{Denormalised counters}: \texttt{Project.proposalCount} and \texttt{Project.viewCount} are maintained as denormalised integers for fast read-path performance rather than computing counts via aggregation on every request.
\end{enumerate}
